\pdfbookmark[0]{Magnetic Observatory. 1834}{magnetisches_observatorium_1834}
\begin{center}
\hrule\vspace{2mm}
Göttingische gelehrte Anzeigen. 1834. 9 August.
\vspace{2mm}\hrule\vspace{5mm}
\end{center}
Thanks to the favour of our government, we have a new institution devoted to an important branch of the natural sciences, an observatory built specifically
\begin{center}\itshape for magnetic observations and measurements\normalfont.\end{center}

Although construction was completed last autumn and the interior arrangements have been far enough along since the beginning of this year for daily observations to have been made since the early months, we have so far held back from reporting on it in these pages. We wanted to be able to include some results of the observations. We previously described in detail in these pages [Notice of \textit{Intensitas v.\,m.}] the magnetic instruments built on new principles and installed in the astronomical observatory here in 1832. The precision they can
\SourceLeaf{526}achieve is clear enough from what we said there. But to attain that precision \textit{in full}, the apparatus had to be built on a larger scale, and to obtain results entirely free from outside influences, a separate iron-free building was indispensable.

The magnetic observatory stands on an open site about a hundred paces west of the astronomical observatory. It is a precisely aligned oblong rectangle, 32 Paris feet long and 15 feet wide, with two projections on the longer sides. The western projection forms the entrance and also serves as an extension of the main hall for certain observations. The eastern projection is entirely separate from the main hall and provides a place for the astronomical observatory’s night watchman to stay. Throughout the building, without exception, everything for which iron is ordinarily used, including locks, door hinges, window fittings and nails, is made of copper. Every effort has been made to keep out draughts. The hall is a little over 10 feet high.

The magnetic apparatus is essentially the same as the instruments mentioned above, so we shall describe only the differences. The magnetic bar is made of cast steel from Uslar, which is excellently suited to magnetic experiments. Different bars are used in turn from time to time. They are all nearly the same size, namely 610 millimetres long, 37 wide and 10 thick. The weight is about four pounds. The mirror is 75 millimetres wide and 50 high. The bar hangs from the centre of the hall ceiling on an untwisted 200-strand silk thread 7 feet long. The torsion circle, however, is at the lower end of the thread instead of the upper end as before, and is joined to the cradle carrying the bar so that it can turn. As Court Councillor \textsc{Gauss} has already noted in his treatise (\textit{Intensitas vis magneticae terrestris}, Article 9), silk suspension threads have the great advantage over metal ones that their torsional force is very small. With the present suspension thread it is only one nine-hundredth of the horizontal directing force of the magnetic bar, whereas the torsional force of a metal thread with the same load-bearing capacity would be about ten times as great. Silk threads, on the other hand, especially when their load-bearing capacity is not much greater than the weight hanging from them, have the drawback that they lengthen considerably in the first weeks or when the load is substantially increased. Here this drawback is remedied by the ingenious ceiling-mounted suspension apparatus
\SourceLeaf{527}devised by Professor \textsc{Weber}. It lets the thread be wound up again as much as necessary without moving it from its place. The apparatus itself can just as easily be shifted along the ceiling when the change in magnetic declination makes this necessary over time. For the present, the theodolite stands on a very solidly built wooden tripod over its own stone foundation. From that position one of the town towers, whose azimuth has been determined with the utmost precision, can be seen through the north window. A fine vertical line on the opposite north wall is the only reference mark used to check that the theodolite has not shifted. The scale in ordinary use is 4 feet long and graduated in millimetres. For some observations it is replaced by one two metres long. One scale division corresponds to $21''3$. Until now the scale has been lit by large wax candles for night observations. In future \textsc{Argand} lamps will be used.

One of the main uses of the apparatus is to determine magnetic declination precisely and to measure how it varies at different hours of the day, in different months and in different years. Readings are recorded twice every day at fixed hours. The hours chosen were 8 in the morning and 1 in the afternoon. When the daily variations follow their regular course, the smallest and largest declinations coincide approximately with these times, at least in the early months of the year. For important reasons that cannot be set out here, recording the readings at exactly the same times on every occasion had to be given unconditional preference over waiting each time for the minimum and maximum. The records began on 1 January. At first, however, a weaker suspension thread was used. It gradually lengthened and had to be wound up frequently, causing appreciable changes in the torsion zero point that initially received too little attention. It was therefore thought better to exclude the first two and a half months. The mean values obtained since then for the westward declination of the magnetic needle were as follows.
\SourceLeaf{528}
\begin{center}
\begin{tabular}{l|rrr|rrr}
 & \multicolumn{3}{c|}{8 a.m.} & \multicolumn{3}{c}{1 p.m.}\\\hline
March, second half & $18^0$ & $38'$ & $16''0$ & $18^0$ & $46'$ & $40''4$\\
April & & $36$ & $6.9$ & & $47$ & $3.8$\\
May & & $36$ & $28.2$ & & $47$ & $15.4$\\
June & & $37$ & $40.7$ & & $47$ & $59.5$\\
July & & $37$ & $57.5$ & & $48$ & $19.0$
\end{tabular}
\end{center}

On certain fixed days of the year, changes in declination are also observed without interruption for 44 hours at short intervals. The days chosen are the same ones that Mr von \textsc{Humboldt} fixed several years ago and on which similar readings are already being recorded by arrangement in many places, some very far away, with \textsc{Gambey} instruments. So far the observations have been made here three times, on 20 and 21 March, 4 and 5 May, and 21 and 22 June. In addition to Court Councillor \textsc{Gauss}, the participants were Professors \textsc{Weber} and \textsc{Ulrich}, Dr \textsc{Weber}, Dr \textsc{Goldschmidt}, Dr \textsc{Listing}, and Messrs \textsc{Sartorius}, \textsc{Deahna} and \textsc{Wilh.\ Gauss}. These observations have two aims. The first is to learn ever more fully about the regular course of the variations. The second is to investigate, by comparing simultaneous observations at different places, the nature of the extraordinary anomalies that so often interrupt that course and can sometimes be remarkably large, especially during displays of the northern lights. Here readings were recorded from 20-minute mark to 20-minute mark in March, and sometimes at half those intervals, from 10-minute mark to 10-minute mark in May, and sometimes twice as frequently, and from 5-minute mark to 5-minute mark throughout June. A couple of strikingly large anomalies were seen here on the night of 20 to 21 March. Very considerable and numerous ones were seen on the nights of 4 and 5 May. Some were distinctly noticeable on 21 June, although they were not large, while the course was exceedingly regular throughout 22 June. As for the corresponding observations that, as already mentioned, were prompted by Mr von \textsc{Humboldt}, the only ones we have heard of so far are those made in Berlin on 20 and 21 March. They were recorded only once an hour and could therefore yield no particular results, although they did show a trace of the anomalies observed and followed in Göttingen. Mr \textsc{Sartorius}, on the other hand, used a smaller instrument built on the same principles as the one
\SourceLeaf{529}here to make very complete corresponding observations on 4 and 5 May at an estate in Bavaria, a few Meilen south of Meiningen. These show a truly remarkable agreement with the large anomalies observed here, in their timing, magnitude and changes. In the graphs one would almost take one for a copy of the other, with all the bizarre shapes formed by the anomalies. The result on 21 and 22 June was just as fine. Corresponding observations were made in Berlin for the first time with an instrument like ours, though smaller, by Professor \textsc{Encke} with the help of Messrs \textsc{Poggendorff}, \textsc{Mädler} and \textsc{Wolfers}. There too the only anomalies were the ones observed here, reproduced almost exactly. The observations Mr \textsc{Sartorius} made on that occasion at Frankfurt am Main showed just the same thing. These results can already be regarded as a fine outcome of the coordinated observations. They show with the greatest clarity that the smaller and larger anomalies of the magnetic needle, which sometimes change within fairly short intervals, must have powerful causes whose effects reach far beyond any one locality. This had been noted before for very large irregularities associated with the northern lights. As participation in these coordinated observations continues to spread in future, using instruments as convenient as they are precise, for which the prospects are already excellent, we are sure to gain a much fuller understanding of these most remarkable and puzzling phenomena.

Observations of this kind are also frequently made here outside the fixed times, and quite striking anomalies have occasionally turned up. For example, on 14 January in the evening between 8 and 9 o’clock, the declination decreased by 13 minutes of arc within a single quarter of an hour with the greatest regularity, and then gradually returned to its previous position. Such observations can yield no further results, however, since corresponding observations can very rarely be expected without prior arrangement.

From time to time the absolute intensity of the Earth’s magnetism will also be determined again at the magnetic observatory here. Since various arrangements first had to be made to carry out this operation with the utmost precision, the first determination could only be made
\SourceLeaf{530}in July. Three determinations with different bars gave
\begin{center}
\begin{tabular}{rllllllr}
17 July & . & . & . & . & . & & $1.7743$\\
20. \enspace „ & . & . & . & . & . & & $1.7740$\\
21. \enspace „ & . & . & . & . & . & & $1.7761$
\end{tabular}
\end{center}
as the value of the horizontal force. The units used are the second of time, the millimetre and the milligram, just as in the earlier determinations with smaller bars. One will note with pleasure how little those results differ from the present ones.

As with the earlier apparatus installed in the astronomical observatory, provision has now been made for electromagnetic experiments and measurements with the present apparatus in the magnetic observatory. The suspended magnetic bar is surrounded by a multiplier coil of 200 turns, designed to allow the use of wire without a covering. The wire is 1100 feet long. By means of a very simply built commutator the observer can reverse the direction of the galvanic current or cut it off altogether at any moment without taking their eye from the telescope.

Here we must mention an impressive installation, so far unique of its kind, closely connected with the arrangements described. We owe it to our Professor \textsc{Weber}. Last year he had already run a double wire connection from the physics cabinet over the houses in town to the astronomical observatory. It has now been extended from there to the magnetic observatory. This makes a large galvanic circuit in which the galvanic current has to travel through nearly nine thousand feet of wire, counting the multiplier coils at both ends. The wire in the circuit is mostly copper wire sold as No. 3, of which one metre weighs eight grams. The wire in the multiplier coil in the magnetic observatory is silver-plated copper wire No. 14, of which there are 2,6 metres to a gram. The installation is well suited to providing opportunities for a great many of the most interesting experiments. One cannot help admiring how putting in a single pair of plates at the far end
\SourceLeaf{531}sets the magnetic bar in motion at once, with a deflection that rises to well over a thousand scale divisions. Even more striking, at least at first, is that a very small pair of plates, for instance a single inch in diameter, using only well water or even distilled water, produces an effect not much smaller than that of a very large pair of plates with strong acid. Yet on closer consideration this is entirely as it should be and only confirms anew the beautiful theory first proposed by \textsc{Ohm}. Increasing the number of pairs of plates, by contrast, increases the effect, and almost in proportion to their number. The ease and certainty with which the commutator controls the direction of the current and the resulting motion of the needle had already prompted trials last year of its use for telegraphic signalling. These were entirely successful even with whole words and short phrases. There is no doubt that it would be possible to establish a direct telegraphic connection in a similar way between two places a considerable number of Meilen apart. But this is naturally no place to develop ideas on the subject any further.

\bigskip
